\documentclass[11pt,a4paper,oneside]{book}
\usepackage[a4paper,margin=2.4cm]{geometry}
\usepackage{amsmath,amssymb}
\usepackage{tcolorbox}
\tcbuselibrary{skins,breakable}
\usepackage{enumitem}
\usepackage{titlesec}
\usepackage{microtype}
\usepackage[hidelinks]{hyperref}

\definecolor{defblue}{RGB}{25,85,150}
\definecolor{exgreen}{RGB}{30,120,80}
\definecolor{warnred}{RGB}{170,50,40}

\titleformat{\chapter}[hang]{\normalfont\huge\bfseries\color{defblue}}{\thechapter\hspace{0.6em}\textcolor{gray!60}{|}\hspace{0.6em}}{0pt}{}
\titlespacing*{\chapter}{0pt}{10pt}{25pt}

% Boxes share one counter numbered within chapters.
% Use \begin{definition}{Title}{key} and refer to it with \ref{def:key}.
\newcounter{boxnum}[chapter]
\renewcommand{\theboxnum}{\thechapter.\arabic{boxnum}}
\newtcolorbox{defbox}[1]{enhanced,breakable,colback=defblue!5,colframe=defblue,fonttitle=\bfseries,sharp corners=south,title={Definition \theboxnum: #1}}
\newtcolorbox{exbox}[1]{enhanced,breakable,colback=white,colframe=exgreen,coltitle=exgreen,colbacktitle=exgreen!10,fonttitle=\bfseries,boxrule=0.5pt,attach boxed title to top left={xshift=4mm,yshift=-2mm},boxed title style={colframe=exgreen},top=4mm,title={Example \theboxnum: #1}}
\newenvironment{definition}[2]{\refstepcounter{boxnum}\label{def:#2}\begin{defbox}{#1}}{\end{defbox}}
\newenvironment{example}[2]{\refstepcounter{boxnum}\label{ex:#2}\begin{exbox}{#1}}{\end{exbox}}
\newtcolorbox{pitfall}{enhanced,colback=warnred!5,colframe=warnred!5,borderline west={3pt}{0pt}{warnred},sharp corners,fontupper=\small,before upper={\textbf{\color{warnred}Common mistake.}\ }}
\newtcolorbox{summarybox}{enhanced,colback=gray!8,colframe=gray!8,sharp corners,title={Key points},coltitle=black,fonttitle=\bfseries,colbacktitle=gray!20}

% Exercises get their own counter that restarts each chapter.
\newcounter{exercise}[chapter]
\renewcommand{\theexercise}{\thechapter.\arabic{exercise}}
\newcommand{\exercise}{\refstepcounter{exercise}\par\medskip\noindent\textbf{Exercise \theexercise.}\ }

\title{\color{defblue}\Huge\bfseries Introduction to Probability\\[0.5em]\Large\normalfont Lecture Notes}
\author{Lecturer Name\\ \small Course code, Example University}
\date{Autumn term 2026}

\begin{document}

\frontmatter
\maketitle
\chapter{Preface}
Explain who the notes are for, what background is assumed and how to use the boxes: blue for definitions, green for worked examples and red for common mistakes.

\tableofcontents

\mainmatter

\chapter{Sample spaces and events}

\section{Outcomes and events}
Start each section with a short motivating paragraph before any formal definition.

\begin{definition}{Sample space}{sample}
The \emph{sample space} $\Omega$ of an experiment is the set of all possible outcomes. An \emph{event} is a subset $A \subseteq \Omega$.
\end{definition}

\begin{example}{Rolling a die}{die}
For one roll of a fair die, $\Omega = \{1,2,3,4,5,6\}$. The event ``the roll is even'' is $A = \{2,4,6\}$, so $P(A) = 3/6 = 1/2$.
\end{example}

As Definition~\ref{def:sample} shows, events are just sets, so set operations apply.

\begin{pitfall}
Confusing an outcome $\omega$ with the event $\{\omega\}$. Probabilities are assigned to events.
\end{pitfall}

\section{Axioms of probability}
\begin{definition}{Probability measure}{measure}
A function $P$ on events is a probability measure if $P(A) \ge 0$, $P(\Omega) = 1$ and
\[
  P\Bigl(\bigcup_{i=1}^{\infty} A_i\Bigr) = \sum_{i=1}^{\infty} P(A_i)
\]
for pairwise disjoint events $A_1, A_2, \dots$
\end{definition}

\begin{summarybox}
\begin{itemize}[nosep]
  \item A sample space lists every outcome.
  \item Events are subsets; probabilities attach to events.
  \item Three axioms determine every other rule.
\end{itemize}
\end{summarybox}

\section*{Exercises}
\exercise Write down the sample space for tossing two coins.
\exercise Show that $P(\emptyset) = 0$ using the axioms.
\exercise\label{exr:union} Prove that $P(A \cup B) = P(A) + P(B) - P(A \cap B)$.

\chapter{Conditional probability}
\section{Definition}
\begin{definition}{Conditional probability}{cond}
For $P(B) > 0$, the probability of $A$ given $B$ is $P(A \mid B) = P(A \cap B)/P(B)$.
\end{definition}

\begin{example}{Two cards}{cards}
Draw two cards without replacement. The chance the second is an ace given the first was an ace is $3/51$.
\end{example}

\section*{Exercises}
\exercise Use Exercise~\ref{exr:union} to find $P(A \cup B)$ when $A$ and $B$ are independent.

\appendix
\chapter{Solutions to selected exercises}
Give short solutions or hints here.

\end{document}
